\documentclass[11pt]{article}
\usepackage[a4paper,margin=25mm]{geometry}
\usepackage{amsmath,amssymb,amsthm}
\usepackage[hidelinks]{hyperref}
\newtheorem{theorem}{Theorem}
\begin{document}
\title{A rank-one completion counterexample to Conjecture 2835}
\author{Chengzhe Gao}\date{}\maketitle
The asserted criterion says that a rank-$r$ matrix is uniquely determined
by the entries on a graph support $G$ if and only if $G$ contains an
$r$-regular supported subgraph. The sufficiency direction is false already
for a $2\times2$ rank-one matrix.

Take
\[
 A=\begin{pmatrix}1&1\\1&1\end{pmatrix},\qquad
 B=\begin{pmatrix}1&-1\\-1&1\end{pmatrix},\qquad
 G=\{(1,1),(2,2)\}.
\]
Here $G$ is the bipartite observation graph, with two row vertices and two
column vertices. Its two edges form a perfect matching, so $G$ itself is
$1$-regular. In particular it contains a $1$-regular subgraph whether or
not ``subgraph'' is required to be spanning.

\begin{theorem}
The entries on $G$ do not uniquely recover $A$ among rank-one matrices.
\end{theorem}
\begin{proof}
Both matrices have determinant zero and a nonzero entry, hence rank exactly
one. More explicitly,
$A=(1,1)^{\mathsf T}(1,1)$ and
$B=(1,-1)^{\mathsf T}(1,-1)$.
Their two observed diagonal entries agree. Their $(1,2)$ entries differ,
so they are distinct completions of the same observations, each of the
specified rank. Since $G$ itself is $1$-regular, the conjectured sufficient
condition holds while unique recovery fails.
\end{proof}

The entries are rational, and the same matrices have rank one over the
real and complex fields. The example does not impose an unintended
spanning condition on the source. It also does not use a counting heuristic
as a substitute for constructing two completions. Failure of this necessary
part of the conjecture suffices, independently of the later observation-count
claim.

\paragraph{Formal verification.}
\texttt{Matrix} is the full space of $2\times2$ matrices over
\texttt{Std.Internal.Rat}. The function \texttt{rank} is the determinantal
rank: two if the sole full-size determinant is nonzero, zero if every
one-by-one minor is zero, and one otherwise. This is exactly the standard
rank definition by the largest nonvanishing minor in dimension two.
The proof also checks every selected two-by-two minor and an explicit
nonzero one-by-one minor for each matrix, as well as both outer-product
descriptions.

The support is an actual Boolean bipartite graph. Row and column degrees
count all neighbors, and the regular-subgraph predicate permits arbitrary
selected nonempty row and column vertex sets, with edges restricted to
those sets and to $G$. The witness selects every vertex and $G$ itself.
Unique recoverability quantifies over \emph{all} rational matrices of the
same rank, not a finite candidate table. The program constructs $B$ as a
counterexample to that universal quantifier. The final theorem negates the
rank-one, two-by-two instance of the source's iff criterion. Lean 4.19.0
checks this with \texttt{Std}, no unproved declarations, and no native
evaluation tactic.

\end{document}
